\section{Online Prefix-Free Encoding}
\label{sec:online-prefix-free-encoding}

\subsection{The Online Universal-Coding Problem}

The paper's second main result concerns a binary stream whose final length
is not known when encoding begins. A valid encoding must be prefix-free, so
the decoder can identify the message boundary without receiving the length
separately. An offline Elias code can place an already known length before
the payload, but an online encoder cannot emit such a header before learning
the end of the stream. Adding an integral end marker to every fixed-size
block is online, but incurs linear redundancy.

The construction combines slowly growing blocks, an early termination
symbol, and the local mapping of Section~\ref{sec:information-carriers}. It
establishes a universal code of length
\[
    n+\log_2 n+O(\log\log n)
\]
for an \(n\)-bit input, using \(O(\log n)\) bits of working memory and
constant time per word-level block operation. This develops the simplified
fixed-block intuition of SOLE from Section~\ref{sec:sole} into a code for an
unknown-length binary stream.

\subsection{Slowly Growing Blocks}

Block sizes are determined online from the amount of input already seen.
When the observed prefix has the current scale denoted by \(n\) in the
paper's local description, the next complete block has
\[
    2k_i=2\left\lceil\log_2 n\right\rceil
\]
bits. This use of \(n\) does not assume knowledge of the eventual stream
length. The complete block consists of two \(k_i\)-bit components and is
interpreted as
\[
    x_i\in[2^{k_i}]^2.
\]
Thus \(k_i\) is the component length, whereas the full block length is
\(2k_i\).

The universe of a complete block has size \(2^{2k_i}\). Since the
information-carrier loss on a universe of size \(X\) is
\(O(1/\sqrt X)\), this choice gives a local loss of
\[
    O\left(\frac{1}{\sqrt{2^{2k_i}}}\right)
    =O(2^{-k_i}).
\]
As the stream grows, later rounding losses consequently become smaller.

\subsection{Handling the Final Partial Block}

Let \(N\) be the final number of blocks. The first \(N-1\) blocks are
complete, whereas block \(N\) may contain between one and \(2k_N\) bits.
The main construction encodes the complete prefix, after which it appends an
Elias encoding for the last partial block. This special encoding supplies
only the boundary information for a block of logarithmic size; it does not
place the unknown total length before the entire stream. Its overhead is
\[
    O(\log k_N)=O(\log\log n).
\]

\subsection{Pass One: Embedding an Early EOF Symbol}

Pass one enlarges the alphabet of each complete block from
\([2^{k_i}]^2\) to
\[
    [2^{k_i}]
    \times
    \left([2^{k_i}]\cup\{\mathrm{eof}\}\right).
\]
If block \(N-2\) originally equals
\[
    (a,b)\in[2^{k_{N-2}}]^2,
\]
the encoder replaces it by \((a,\mathrm{eof})\). The displaced component
\(b\) is not discarded: it is appended in plaintext as part of the special
tail.

The encoder maintains a three-block lag using an \(O(\log n)\)-bit buffer.
When a block first arrives, the encoder cannot know whether it will
eventually be block \(N-2\). It therefore delays the last constant number of
blocks and chooses the termination position only after the input ends. No
previously emitted output is rewritten.

The position \(N-2\), rather than the last block, is forced by decoding
locality. Recovering an input block requires information committed to the
next output block. Moving termination sufficiently far forward ensures that
the decoder learns it before exhausting the ordinary encoded stream.

\subsection{Redundancy of Alphabet Enlargement}

The original universe of block \(i\) has \(2^{2k_i}\) states, while the
enlarged universe has \(2^{k_i}(2^{k_i}+1)\) states. Its information cost is
\[
\begin{aligned}
    &\log_2\left(2^{k_i}(2^{k_i}+1)\right)-2k_i\\
    &\qquad=\log_2(1+2^{-k_i})
    =O(2^{-k_i}).
\end{aligned}
\]

The sum is analyzed by block scale. For a fixed \(K\), the slowly growing
partition contains
\[
    O\left(\frac{2^K}{K}\right)
\]
blocks satisfying \(k_i=K\). Since each contributes \(O(2^{-K})\), the
entire scale contributes
\[
    O\left(\frac{2^K}{K}\right)O(2^{-K})
    =O\left(\frac{1}{K}\right).
\]
Summing this harmonic bound over \(1\leq K\leq\log_2n\) gives
\[
    \sum_{K=1}^{\log_2 n}O\left(\frac{1}{K}\right)
    =O(\log\log n).
\]

\subsection{Pass Two: Conversion through Information Carriers}

Pass two converts the enlarged alphabets to binary. As in the sequential
carrier chain preceding the tree construction of
Section~\ref{sec:vector-representation}, it applies the information-carrier
mapping in order:
\[
    (x_i,s_{i-1})\longmapsto(m_i,s_i).
\]
Here \(x_i\) is the current enlarged-alphabet block, \(s_{i-1}\) is the old
spill, \(m_i\) is binary memory content that can be emitted, and \(s_i\) is
the new spill passed to the next block. The asymmetric inverse guaranteed in
Section~\ref{sec:information-carriers} commits the old spill to \(m_i\),
which is precisely the locality needed by the decoder.

The carrier domain at step \(i\) has size \(O(2^{2k_i})\), so the lemma adds
\(O(2^{-k_i})\) redundancy. The same scale decomposition used above bounds
the cumulative carrier redundancy by \(O(\log\log n)\). The proof of the
lemma need not be repeated: its bounded spill and constant number of word
arithmetic operations apply directly to these growing domains.

\subsection{Why the Decoder Stops in Time}

The termination symbol occurs in input block \(N-2\). After reading output
block \(N-1\), the local inverse mapping has recovered the spill needed to
decode input block \(N-2\), and hence reveals its second component
\(\mathrm{eof}\). The decoder therefore sees termination before it would
need to read beyond the ordinary encoded stream.

Upon detecting \(\mathrm{eof}\), the decoder stops the sequential carrier
procedure and switches to the special tail. It decodes the final partial
block and reads the displaced component \(b\), restores \(b\) as the second
half of block \(N-2\), and outputs the remaining original bits. Because the
switch occurs immediately after output block \(N-1\), the special tail is
not mistaken for additional ordinary carrier blocks.

\subsection{Total Encoding Length}

Replacing the second component of block \(N-2\) by termination and moving
that component to the tail costs at most
\(k_{N-2}\leq\log_2n\) bits. This is one termination mechanism, not two
independent charges for the marker and the displaced data. Alphabet
enlargement and the sequential carrier mappings each add
\(O(\log\log n)\), and the final partial block adds another
\(O(\log\log n)\). Consequently,
\[
\begin{aligned}
    \text{encoding length}
    &=n+\underbrace{\log_2n}_{\text{early termination}}
      +\underbrace{O(\log\log n)}_{\text{alphabet enlargement}}\\
    &\quad+\underbrace{O(\log\log n)}_{\text{carrier mappings}}
      +\underbrace{O(\log\log n)}_{\text{partial block}}\\
    &=n+\log_2n+O(\log\log n).
\end{aligned}
\]

\subsection{Memory and Running Time}

The encoder and decoder retain a constant number of \(O(\log n)\)-bit
blocks, one bounded spill, and a constant number of counters and block
parameters. Their working memory is therefore \(O(\log n)\) bits.

Each pass performs a constant number of additions, multiplications,
divisions, and quotient--remainder operations per block. Under the paper's
Word RAM assumptions this is \(O(1)\) time per block, or per word-level
operation. The special ending processes only a constant number of blocks;
the total processing time over a stream remains linear in the number of
processed blocks, rather than constant for the entire input.

\subsection{Significance of the Construction}

SOLE demonstrates with fixed blocks how short and long ranges can carry a
non-binary spill. The general carrier lemma removes that simplifying
structure and supports the varying alphabets required here. Slowly growing
blocks make local rounding losses summable, while early termination aligns
the code boundary with the lemma's local inverse. The same primitive that
gave exact-space vector representation thus also yields the paper's second
main result: a low-memory online code whose redundancy is within
\(O(\log\log n)\) of the Elias-scale bound.
